% GENERATED FILE. Edit canonical lessons, then run npm run generate:books.
% canonical-source-hash: fnv1a64:bfb4246d07306424
% canonical-lessons: ES-C680-pasteleria, ES-C680-papeleria, ES-C680-peluqueria, ES-C680-lavanderia, ES-C680-heladeria

\chapter{Tiendas del barrio}
\label{ch:es-tiendas-del-barrio}
\hlchaptermodality{680}

\begin{chapteropening}
\textbf{By the end of this chapter:} \emph{I can find the cake shop, the stationer’s, the hairdresser, the launderette and the ice-cream shop.}

Complete the last lesson of chapter 680: I can find the cake shop, the stationer’s, the hairdresser, the launderette and the ice-cream shop.
\end{chapteropening}

\section[la pastelería]{la pastelería — the cake shop, the pastry shop}

\label{lesson:ES-C680-pasteleria}



\begin{quote}
Before the new one: say the Spanish for the stamp, then the Spanish for the toilet.
\end{quote}

\subsection*{You'll want to know: la pastelería}
\textbf{la pastelería} — \textquotedblleft{}the cake shop, the pastry shop\textquotedblright{}.

\begin{quote}
Compramos un pastel en la pastelería.
\end{quote}

\begin{quote}
\emph{We bought a cake at the cake shop.}
\end{quote}

The baker of cakes is \emph{el pastelero}.

\begin{cousinweb}
\emph{pastelería} is built on \emph{pastel} \textquotedblleft{}cake\textquotedblright{} with \emph{-ería} (EtymDB); \emph{pastel} is usually traced to Late Latin \emph{pasta} \textquotedblleft{}dough\textquotedblright{}.
\end{cousinweb}

\begin{grammarlens}[title={What you've built}]
One more word for telling, explaining and describing at B1.
\end{grammarlens}

\subsection*{Your turn}
\begin{itemize}
\raggedright
  \item \emph{Say it:} \emph{la pastelería}
  \item \emph{Say it:} \emph{la pastelería}, once more
  \item \emph{Say it:} a sentence of your own with \emph{la pastelería}
  \item \emph{Recall:} say the Spanish for the stamp, then the Spanish for the toilet, then say \emph{la pastelería} again
\end{itemize}

\begin{tcolorbox}[breakable,colback=teal!4,colframe=teal!35!black,title={Before you move on}]
What does la pastelería mean? (\textquotedblleft{}The cake shop, the pastry shop\textquotedblright{}.) Say it once more.
\end{tcolorbox}
\section[la papelería]{la papelería — the stationer’s}

\label{lesson:ES-C680-papeleria}



\begin{quote}
Before the new one: say the Spanish for the toilet, then the Spanish for the cake shop.
\end{quote}

\subsection*{You'll want to know: la papelería}
\textbf{la papelería} — \textquotedblleft{}the stationer’s\textquotedblright{}.

\begin{quote}
Venden cuadernos y bolígrafos en la papelería.
\end{quote}

\begin{quote}
\emph{They sell notebooks and pens at the stationer’s.}
\end{quote}

For pens, paper and envelopes.

\begin{cousinweb}
\emph{papelería} is built on \emph{papel} \textquotedblleft{}paper\textquotedblright{}, which came from Catalan \emph{paper}, from Latin \emph{papyrus} (EtymDB).
\end{cousinweb}

\begin{grammarlens}[title={What you've built}]
One more word for telling, explaining and describing at B1.
\end{grammarlens}

\subsection*{Your turn}
\begin{itemize}
\raggedright
  \item \emph{Say it:} \emph{la papelería}
  \item \emph{Say it:} \emph{la papelería}, once more
  \item \emph{Say it:} a sentence of your own with \emph{la papelería}
  \item \emph{Recall:} say the Spanish for the toilet, then the Spanish for the cake shop, then say \emph{la papelería} again
\end{itemize}

\begin{tcolorbox}[breakable,colback=teal!4,colframe=teal!35!black,title={Before you move on}]
What does la papelería mean? (\textquotedblleft{}The stationer’s\textquotedblright{}.) Say it once more.
\end{tcolorbox}
\section[la peluquería]{la peluquería — the hairdresser’s}

\label{lesson:ES-C680-peluqueria}



\begin{quote}
Before the new one: say the Spanish for the cake shop, then the Spanish for the stationer’s.
\end{quote}

\subsection*{You'll want to know: la peluquería}
\textbf{la peluquería} — \textquotedblleft{}the hairdresser’s\textquotedblright{}.

\begin{quote}
La peluquería cierra los lunes.
\end{quote}

\begin{quote}
\emph{The hairdresser’s is closed on Mondays.}
\end{quote}

The hairdresser is \emph{el peluquero}, \emph{la peluquera}.

\begin{cousinweb}
\emph{peluquería} is built on \emph{peluca} \textquotedblleft{}wig\textquotedblright{} (EtymDB), which came from French \emph{perruque} and Italian before it.
\end{cousinweb}

\begin{grammarlens}[title={What you've built}]
One more word for telling, explaining and describing at B1.
\end{grammarlens}

\subsection*{Your turn}
\begin{itemize}
\raggedright
  \item \emph{Say it:} \emph{la peluquería}
  \item \emph{Say it:} \emph{la peluquería}, once more
  \item \emph{Say it:} a sentence of your own with \emph{la peluquería}
  \item \emph{Recall:} say the Spanish for the cake shop, then the Spanish for the stationer’s, then say \emph{la peluquería} again
\end{itemize}

\begin{tcolorbox}[breakable,colback=teal!4,colframe=teal!35!black,title={Before you move on}]
What does la peluquería mean? (\textquotedblleft{}The hairdresser’s\textquotedblright{}.) Say it once more.
\end{tcolorbox}
\section[la lavandería]{la lavandería — the launderette, the laundry}

\label{lesson:ES-C680-lavanderia}



\begin{quote}
Before the new one: say the Spanish for the stationer’s, then the Spanish for the hairdresser’s.
\end{quote}

\subsection*{You'll want to know: la lavandería}
\textbf{la lavandería} — \textquotedblleft{}the launderette, the laundry\textquotedblright{}.

\begin{quote}
En el barrio hay una lavandería.
\end{quote}

\begin{quote}
\emph{There is a launderette in the neighbourhood.}
\end{quote}

\emph{Lavandería automática} — self-service launderette.

\begin{cousinweb}
\emph{lavandería} is built on \emph{lavandero} \textquotedblleft{}washerman\textquotedblright{}, from \emph{lavar} \textquotedblleft{}to wash\textquotedblright{}, which is inherited from Latin \emph{lavāre} (EtymDB).
\end{cousinweb}

\begin{grammarlens}[title={What you've built}]
One more word for telling, explaining and describing at B1.
\end{grammarlens}

\subsection*{Your turn}
\begin{itemize}
\raggedright
  \item \emph{Say it:} \emph{la lavandería}
  \item \emph{Say it:} \emph{la lavandería}, once more
  \item \emph{Say it:} a sentence of your own with \emph{la lavandería}
  \item \emph{Recall:} say the Spanish for the stationer’s, then the Spanish for the hairdresser’s, then say \emph{la lavandería} again
\end{itemize}

\begin{tcolorbox}[breakable,colback=teal!4,colframe=teal!35!black,title={Before you move on}]
What does la lavandería mean? (\textquotedblleft{}The launderette, the laundry\textquotedblright{}.) Say it once more.
\end{tcolorbox}
\section[la heladería]{la heladería — the ice-cream shop}

\label{lesson:ES-C680-heladeria}



\begin{quote}
Before the new one: say the Spanish for the hairdresser’s, then the Spanish for the launderette.
\end{quote}

\subsection*{You'll want to know: la heladería}
\textbf{la heladería} — \textquotedblleft{}the ice-cream shop\textquotedblright{}.

\begin{quote}
Después de cenar vamos a la heladería.
\end{quote}

\begin{quote}
\emph{After dinner we go to the ice-cream shop.}
\end{quote}

Built on \emph{el helado}, ice cream.

\begin{cousinweb}
\emph{heladería} is built on \emph{helado} \textquotedblleft{}ice cream\textquotedblright{}, the past form of a verb inherited from Latin \emph{gelāre} \textquotedblleft{}to freeze\textquotedblright{}.
\end{cousinweb}

\begin{grammarlens}[title={What you've built}]
One more word for telling, explaining and describing at B1.
\end{grammarlens}

\subsection*{Your turn}
\begin{itemize}
\raggedright
  \item \emph{Say it:} \emph{la heladería}
  \item \emph{Say it:} \emph{la heladería}, once more
  \item \emph{Say it:} a sentence of your own with \emph{la heladería}
  \item \emph{Recall:} say the Spanish for the hairdresser’s, then the Spanish for the launderette, then say \emph{la heladería} again
\end{itemize}

\begin{tcolorbox}[breakable,colback=teal!4,colframe=teal!35!black,title={Before you move on}]
What does la heladería mean? (\textquotedblleft{}The ice-cream shop\textquotedblright{}.) Say it once more.
\end{tcolorbox}
